% =====================================================================
% §2 Related work
% =====================================================================

\section{Related work}
\label{sec:relwork}

We position this paper at the intersection of four threads: time-series forecasting with Transformers, foundation models for time series, electricity demand forecasting, and distribution shift in deployed ML.

\paragraph{Transformer architectures for time-series forecasting.} The Transformer \citep{vaswani2017attention} is now the default backbone for sequence modelling, and a sequence of papers has adapted it to time-series. ViT \citep{dosovitskiy2021vit} introduced patch tokenization for images, which we borrow with a CNN trunk for our coarse $450{\times}449{\times}7$ rasters. Earthformer \citep{gao2022earthformer} explored space-time Transformers for Earth-system forecasting and motivated our shared-CNN-then-Transformer design. The Temporal Fusion Transformer \citep{lim2021tft} introduced an encoder-decoder design with variable-selection networks for multi-horizon forecasting; our encoder-decoder ablation in \S\ref{sec:archsearch} is in that lineage. PatchTST \citep{nie2023patchtst} demonstrated channel-independent temporal patching as a strong univariate-Transformer baseline; this informs our treatment of demand history in the baseline's tabular branch.

\paragraph{Foundation models for time-series.} The recent wave of pretrained time-series foundation models \citep{ansari2024chronos,das2024timesfm,rasul2024lagllama,woo2024moirai} promises zero-shot forecasting on previously-unseen series. We use Chronos-Bolt \citep{ansari2024chronos} because its tokenized-residual design and Apache-2.0 license make it deployable on HuggingFace Spaces' free CPU tier. The mini variant ($\sim$21\,M parameters) is small enough to load alongside our 1.75\,M-parameter baseline within the 16\,GB RAM limit; the base variant ($\sim$205\,M parameters) is reported here as a reference upper bound but is not deployed. We do not run TimesFM, Lag-Llama, or Moirai head-to-head, because (a) Chronos-Bolt-mini already achieves a 17\,\% relative improvement over our baseline at acceptable cost, and (b) the focus of this paper is the deployment behaviour, not foundation-model selection.

\paragraph{Electricity demand forecasting.} Short-term load forecasting has a long classical literature. Box--Jenkins ARIMA \citep{box1976time} and exponential smoothing remain hard benchmarks for short horizons; \citet{hippert2001neural} review classical neural baselines (MLPs, recurrent nets) on this task and document the temperature-sensitive structure that any modern model must respect. \citet{hong2016probabilistic} survey probabilistic forecasting methods for the next decade. We do not run head-to-head ARIMA / ETS comparisons in this paper; \citet{hippert2001neural}'s observations on temperature-sensitive demand and weekday/weekend / holiday seasonality remain the design constraints we honor through the calendar one-hot and HRRR weather inputs. Recent deep-learning work on grid-scale forecasting has emphasized either Transformer architectures or weather-aware multi-modal fusion, but to our knowledge no published paper combines these with a frozen foundation-model partner under per-zone weighted late fusion, which is the design we report on here.

\paragraph{Distribution shift in deployed ML.} The dataset-shift literature \citep{quinonero2009dataset,koh2021wilds} formalizes the gap between train-time and test-time distributions. Most benchmarks treat shift as a static property of carefully-constructed test sets; the energy-transition story is closer to a true \emph{temporal} shift on a single deployed model. \citet{koh2021wilds} include weather and load datasets but their shift definitions are domain-curated rather than measured against a continuously-updating live feed. The novelty of Part II is that we observe the shift in the wild, on a model whose training data closes in late 2022, deployed against ISO-NE 2026 with full pipeline-equivalence to the training-time evaluation.
